\chapter{Installation \& deployment}
\label{ch:install}

\section{Prerequisites}

\prog is a single container. It needs four things from around it, and none of them
is optional:

\begin{itemize}
  \item \textbf{A Docker host.} The published image is \code{linux/amd64} only. That
        is a Dockerfile line rather than an architectural position --- the build sets
        the \code{mold} linker for \code{x86\_64-unknown-linux-gnu} alone, and an
        arm64 image would need that extended first.
  \item \textbf{An MQTT broker.} \prog connects to an \emph{external} broker and does
        not embed one. It must be reachable from the container, which in practice
        means the two share a Docker network and the broker is addressed by its
        \emph{container name}, not by a LAN address.
  \item \textbf{Outbound HTTPS to the smart-me cloud.} The API base must be TLS: the
        client refuses a plain-HTTP endpoint outright rather than degrading.
  \item \textbf{A smart-me OAuth2 credential} --- a client id and client secret, from
        a ``Device'' type application in your smart-me account. This is the one
        setting that never descends to disk (section~\ref{sec:install-env}).
\end{itemize}

A reverse proxy is not required, but is assumed throughout: the container publishes
no host port, so without one you must uncomment the \code{ports:} fallback described
in section~\ref{sec:install-traefik}.

\section{Quick start with \texttt{docker compose}}
\label{sec:install-quickstart}

\subsection{The directory}

Copy \code{docker-compose.yml} and \code{.env.example} from the repository into a
directory of their own, and rename the second to \code{.env}:

\begin{lstlisting}[language=bash]
mkdir -p /srv/docker/smartme_mqtt && cd $_
# copy docker-compose.yml and .env.example here
mv .env.example .env
chmod 600 .env
\end{lstlisting}

\subsection{The state directory, and the step everyone forgets}
\label{sec:install-chown}

\begin{warning}
\textbf{The state directory must belong to uid 10002.} The container runs as an
unprivileged user with that fixed uid, chosen distinct from \code{opcgw}'s 10001 so
neither service can rewrite the other's state.

\begin{lstlisting}[language=bash]
mkdir -p ./data && sudo chown -R 10002:10002 ./data
\end{lstlisting}

Docker creates a missing bind-mount source \emph{itself}, owned by root, so skipping
this does not fail loudly --- the bridge comes up, serves its web UI, and cannot
write the configuration you type into it. The screen says so (``the configuration was
NOT written''), which is the only reason this is a nuisance rather than a mystery.

On a Synology NAS there is a second layer: the displayed mode can disagree with the
enforced ACL, so a directory showing \code{777} still refuses the write. If the mode
carries a trailing \code{+}, ACL inheritance is what governs and must be cut on the
folder before the mode bits mean anything.
\end{warning}

What lives in \code{./data}:

\begin{itemize}
  \item \code{config.toml} --- the whole configuration, written by the web UI;
  \item \code{bdseq.toml} --- the Sparkplug session number;
  \item \code{logs/} --- rotated log files, if you enable them.
\end{itemize}

\begin{warning}
\textbf{\code{bdseq.toml} must survive restarts.} A bridge that comes back with a
fresh state directory replays a session number, and a consumer pairing a death
certificate to a birth by \code{bdSeq} can then discard a death belonging to a
session it believes is still live --- so a node that is gone goes on looking alive.
Never delete \code{./data} to ``start clean''.
\end{warning}

\subsection{The environment: exactly three variables}
\label{sec:install-env}

\begin{lstlisting}
SMARTME_CLIENT_ID=...
SMARTME_CLIENT_SECRET=...
\end{lstlisting}

plus \code{SMARTME\_STATE\_DIR}, which the compose file already sets to \code{/data}.
\code{RUST\_LOG} is honoured if present and defaults to \code{info}.

\begin{note}
\textbf{Eleven variables were withdrawn in v0.4.0} (ADR 0023) --- the Sparkplug
identity, the broker, the meter, the publish period, the log settings. \emph{Setting
them now does nothing at all.} If you are upgrading from v0.3.x, see
section~\ref{sec:install-upgrade}.

Both halves of the credential live here together, and that is a decision rather than
an oversight: a credential split across two sources can be rotated by halves, and a
mismatched pair fails as an authentication rejection from the smart-me API --- which
reads as an \emph{outage of the upstream service} rather than as a configuration
fault, and gets diagnosed accordingly, at length.
\end{note}

\subsection{First start}

\begin{lstlisting}[language=bash]
docker compose up -d
docker compose logs -f
\end{lstlisting}

The bridge comes up, says it has no configuration, serves its web UI, and
\textbf{publishes nothing} --- no connection to the broker, no birth. That is the
first run and it is not an error.

Open the UI (section~\ref{sec:install-traefik} for where), fill in the form, and save.
Then \textbf{confirm the mapping on its own screen}: nothing is published until you
do, and the separation is deliberate. A mapping can be well-formed, unique and
complete and still point \code{garage} at the cellar's meter --- and a SCADA host
\emph{persists what it discovers}, so the tag folder outlives the process and has to
be deleted by hand, discarding the alarm and history configuration beneath it. See
section~\ref{sec:confirm}.

Time from a clean state directory to a publishing bridge is about one second for the
container itself; the budget for the whole exercise, including you typing, is fifteen
minutes.

\section{Behind Traefik}
\label{sec:install-traefik}

The compose file joins two external networks and publishes no host port:

\begin{itemize}
  \item \code{scada} --- where the broker is, reached by container name, and the way
        out to the smart-me cloud;
  \item \code{proxy} --- where Traefik looks. \code{traefik.docker.network} names it
        explicitly so Traefik does not have to guess between the two.
\end{itemize}

Both must already exist (\code{docker network create scada proxy}, or whatever your
other services already use).

The router is one host rule at \code{smartme.home.arpa} on the \code{web}
entrypoint. The port named by the load-balancer label must match the bridge's
\code{ui\_port}, which defaults to 8080. If you change that in the UI, change the
label with it --- and note that the bridge treats \code{ui\_port} as needing a
process restart, so the two move together in any case.

\begin{warning}
\textbf{The UI is at \code{http://smartme.home.arpa} --- with no port number.}
Traefik listens on 80; 8080 is the port \emph{inside} the container, which the
compose file deliberately does not publish. That is what keeps an
unauthenticated UI off the LAN.

Adding \code{:8080} does not fail in a way that says so. Traefik opens that port
for its own API, and with \code{api.insecure=false} it answers \textbf{404 page
not found} to everything --- which reads exactly like a bridge that is not
routed, and sends you to look at labels that are correct. If you get a 404, drop
the port before changing anything else.
\end{warning}

\begin{note}
\textbf{There is no authentication on the UI} (ADR 0019): the trust boundary is
Traefik, not this process, and no host port is published, so the container is not
reachable from the LAN except through the proxy.

Inside the bridge, what replaces the missing login is a \emph{same-origin check}
(ADR 0024). Precisely \emph{because} there is no login, any page open in your browser
could otherwise post to the bridge --- a drive-by reconfiguration needs no credential
to steal when there are none. Submissions carrying a foreign \code{Origin} are
refused; requests carrying none at all are allowed, which is what makes the headless
bring-up with \code{curl} work.
\end{note}

\subsection{Without a reverse proxy}

Uncomment the \code{ports:} block, set the Traefik label to
\code{false}, and reach the UI on port 8080 of the host. Understand what you are
doing: the UI has no authentication, so anyone who can reach that port can
reconfigure the bridge and move every meter into a namespace your SCADA host will
keep.

\section{No healthcheck, on purpose}
\label{sec:install-healthcheck}

The compose file defines no \code{healthcheck:} and the image carries no
\code{HEALTHCHECK} line.

\code{/healthz} exists and answers \code{503} for a wedged poll loop --- a loop that
has stopped looping, which is exactly what a restart fixes. Everything else answers
\code{200}, including an honest \code{STALE} and every deliberate silence, because
restarting a container over those destroys the continuity they exist to report.

\begin{warning}
\textbf{Do not add a healthcheck yet.} The endpoint also answers \code{500} from any
handler that panics, its own included --- the guard that catches a panicking page
wraps every route. Wiring that to a container restart would kill the Sparkplug
session for every meter over one broken display page, which is the opposite of what
the endpoint is for. The decision belongs with the healthcheck wiring itself and has
not been taken.
\end{warning}

\section{Upgrading from v0.3.x}
\label{sec:install-upgrade}

v0.4.0 moved the configuration out of the environment and into
\code{data/config.toml}. An upgrade is therefore not only a \code{pull}:

\begin{enumerate}
  \item \textbf{Write down what the old deployment published under} --- \code{group\_id},
        \code{node\_id}, the broker, and each meter's id, device id and serial. They
        are in the old \code{.env}. \textbf{These must be reproduced exactly}: a
        changed group or node id publishes into a namespace nobody has seen, and the
        old one's tag folder stays behind.
  \item Bump the \code{image:} tag, then remove the withdrawn variables from the
        compose file and from \code{.env}. Leaving them costs nothing but reads as
        though they still did something.
  \item Start the bridge. With no \code{config.toml} it comes up unconfigured and
        publishes nothing --- your old namespace is untouched while you work.
  \item Enter the same values in the UI, check each serial against its topic, and
        confirm.
\end{enumerate}

\begin{note}
\textbf{One field cannot be guessed and is worth checking twice.} \code{serial} must
be the serial smart-me \emph{reports} for that \code{device\_id}. The bridge births
the device under the configured serial and routes every reading by the reported one;
before v0.4.0 a mismatch produced a bridge that fetched, judged, ticked and looked
healthy on every surface while publishing nothing at all. It now refuses the meter
and names both serials in the log --- but the fastest way not to meet that refusal is
to let the discovery screen fill the field.
\end{note}
